% Ad Atticum 9.9 (ad-atticum-09-09) --- English translation
% Translated by Claude (Anthropic) from the Latin (Perseus).
% Style guide: STYLE.md.

\ciceroLetterOpener{CICERO ATTICO salutem}

\ciceroSection{1} Three letters of yours I received the day after the Ides; they had been sent, however, on the fourth, the third, and the day before the Ides. Accordingly I shall reply to the oldest first. I agree with you that I should stay on at the Formian villa above all others; also about the Upper Sea --- I shall consider \textit{[Greek: plaboque, text corrupt]}, as I wrote to you before, by what method I can, with his goodwill, refrain from touching any part of the commonwealth. As for your praise of me because I wrote that I am forgetting our friend's earlier acts and faults: I am, indeed, doing so. Nay more, I do not bear in mind that those very things which you recall were done by him against me otherwise than they were. I would have the gratitude for benefits weigh with me so much more than the resentment of wrongs. So let us do as you advise, and pull ourselves together. For I play the philosopher \textit{[Greek: sophisteu\=o]} the moment I run down into the country, and in the running do not stop turning over my propositions \textit{[Greek: theseis]}. But some of them are extremely hard to settle. About the optimates, very well, let it be as you wish; but you know that line, ``Dionysus in --- ''\textit{[Greek: Dionysios en, text broken]}. Titinius's son is with Caesar. As for your seeming, in a way, to fear that your advice may not please me --- nothing else delights me but your counsel and your letters. So, as you propose, do not leave off writing me whatever comes into your mind. Nothing can be more welcome to me.

\ciceroSection{2} I come now to the second letter. You are right not to believe the number of soldiers; Clodia wrote it just half too large. False also about the ships destroyed. Where you praise the consul, I too praise his spirit, but find fault with his judgment; for by their being scattered, that action toward peace which I was meditating has been swept away. So afterwards I sent Demetrius's book on concord back to you and gave it to Philotimus. Nor do I doubt that a ruinous war hangs over us, of which the beginning will be drawn from famine. And yet I grieve that I have no part in this war! In which there will be such force of wickedness that, while it is unspeakable not to feed one's parents, our leaders think the most ancient and most sacred parent of all, the fatherland, should be killed by starvation. And this I fear not by guesswork but as one who has been present at the talk. The whole of this fleet from Alexandria, Colchis, Tyre, Sidon, Aradus, Cyprus, Pamphylia, Lycia, Rhodes, Chios, Byzantium, Lesbos, Smyrna, Miletus, Cos is being prepared to cut off Italy's supplies and to seize the grain-bearing provinces. And how full of anger he will come! and most of all against those who most wanted him safe, as though he had been deserted by those whom he himself deserted. So, for me as I hesitate what it is fitting for me to do, my goodwill toward him brings vast weight; remove that, and it would be better to perish in the fatherland than, by saving the fatherland, to overturn it. About the north of the country, plainly it is so. I fear that Epirus may be harassed; but what place in Greece do you think will not be plundered? For he openly proclaims, and shows to his soldiers, that in the very matter of largesse he will outdo this man here. That, too, is a fine point of yours: when I see him, not to speak too cringingly, and to speak rather with weight. Plainly that is how it must be done. I am thinking of Arpinum, after I have met with him, lest I should happen either to be absent when he comes or to be running this way and that on a very bad road. Bibulus, as you write, I hear has come and gone back the day before the Ides.

\ciceroSection{3} You were waiting, you say in the third letter, for Philotimus. But he set out from me on the Ides. So my letter, to which I had at once written back to that letter of yours, was delivered later than expected. About Domitius, as you write, I take it to be so --- that he is on his estate at Cosa, and that his plans are not known. That most disgraceful and basest of men, who says that consular elections can be held by a praetor, is the same fellow he has always been in public life. So, no doubt, this is what Caesar means in that letter of which I sent you a copy --- when he says he wants to use my advice (come, very well; that is open to all), my influence (an inept word, but I suppose he affects it for the sake of certain votes among the senators), my standing (a consular vote, perhaps). The last item is the sting: ``and my resources of every kind.'' I began to suspect, from your letter, that either this is exactly what it means or not far from it. For it is of great consequence to him that the matter not come to an interregnum. He gets that, if consuls are elected by a praetor. We have it in our books, however, that not only are consuls not lawfully elected by a praetor, but not even praetors --- and that it has never been done; that for consuls the right does not exist, because the greater command may not lawfully be put up by the lesser; for praetors, because they are put up in such a way as to be colleagues of the consuls, whose command is the greater. He will not be far from wanting this decreed by me, and not be content with Galba, Scaevola, Cassius, Antonius. ``Then let the wide earth gape for me'' \textit{[Greek: tote moi chanoi eureia chth\=on]}.

\ciceroSection{4} But you see what a tempest is hanging over us. What senators have crossed over, I shall write to you when I know for certain. About the grain supply, you understand rightly --- it cannot in any way be administered without revenues; and not without cause do you fear both the men around him, who demand everything, and an unspeakable war. Our Trebatius, although, as you write, he hopes for no good, I should still very much like to see. Do urge him to make haste; for it will fall out conveniently if he comes to me before Caesar's arrival. About the Lanuvine place, as soon as I heard Phamea was dead, I hoped --- if only there should be a commonwealth still --- that some one of my friends would buy it; and yet I did not think of you, who are most of all my own. For I knew the yearly rate and the rate-per-acre you are accustomed to seek, and I had seen your land-plan \textit{[Greek: diagramma]} not only at Rome but at Delos. All the same I value the property, though it is charming, at less than it was valued at in Marcellinus's consulship, when I was thinking that those little gardens would be more pleasant to me because of an old house I then had near by, and at less expense than if I rebuilt the Tusculan place. I wanted to give five hundred thousand sesterces. I dealt through \textit{[a friend whose name the manuscript loses]}, that he should let me have it at that, since he had it for sale. He refused. But now I think all those things are flat, because of the scarcity of cash. To me, indeed --- or rather to us both --- it will be most suitable if you buy it; but do not despise the man's strange fancies. The place is very pretty. And yet all those properties seem to me now made over to ruin. I have answered three letters, but I am looking for others; for thus far your letters have been my support. Given on the Liberalia.
